\section{Introduction}\label{sec:intro}
Teams of unmanned aerial vehicles (UAVs) are used for inspection, delivery and monitoring. During a mission, new tasks
may appear, a vehicle may fail, or a deadline may be moved forward. A central planner then repairs the routes of the
whole fleet within a fraction of a second, while some UAVs may already wait for new instructions. Dynamic vehicle
routing and multi-robot task allocation offer many methods for this
repair~\cite{gerkey2004formal,pillac2013review,psaraftis2016dynamic}, and metaheuristics such as adaptive large
neighbourhood search (ALNS)~\cite{ropke2006adaptive} find good plans within such budgets.

Most of these replanners plan for nominal execution. When UAVs meet wind, fly slower than commanded or take longer at a
task, part of the slack that the plan promises is consumed by errors that the planner does not see, and tasks near the
end of a route are served late. The error grows with the time until a task is reached, so a constant time margin is too
large for early tasks and too small for late ones. Stochastic and robust routing protect schedules against such
errors~\cite{errico2016apriori,agra2013robust}, and chance-constrained allocation bounds the probability of
failure~\cite{ponda2012chance}. These models are usually solved by sampling or by exact methods that are too slow for
the time available between two events. A second line of work treats the computation time of an anytime planner as a
decision~\cite{dean1988analysis,russell1991right,zilberstein1996using,svegliato2018meta}. It is not known which of the
two, the planning model or the control of computation, matters more for a replanner that faces uncertain execution.

This paper proposes AnytimeSwarm-CC, a replanner that plans against a chance constraint on the finish time of every
task and keeps the fast heuristics of the nominal model. Its contributions are threefold. The first is a closed-form
chance-constrained planning model. A calibrated finish-time error, whose mean and spread grow linearly with the lead
time, turns a chance constraint on every task into an effective deadline. The margin thus grows along a route, and the
model is solved by the same repair and ALNS as the nominal one, with a negligible overhead per call. A nominal fallback
keeps the tasks that the constraint would give up and leaves the guarantee of all other tasks intact. The second is a
pre-registered comparison on held-out missions with re-implemented baselines. These include stopping rules from the
metareasoning literature~\cite{boddy1994deliberation,zilberstein1996using,svegliato2018meta}, size-scaled and fixed
budgets, a constant time buffer~\cite{agra2013robust} and a conventional event-triggered planner, and ablations
separate the effects of the planning model, the stopping rule and the replanning trigger. The third is a
software-in-the-loop check with PX4 autopilots~\cite{meier2015px4} in Gazebo~\cite{koenig2004gazebo} and
ROS~2~\cite{macenski2022ros2}, after a rescaling that keeps every dimensionless ratio of the decision problem, together
with an estimate of the number of flights that such a check requires.

The results are reported as they came out. In the kinematic tier, the chance-constrained model served \val{KDabsBuf} to
\val{KDabsForty} percentage points more reward than each of six baselines, and ablations traced the gain to the planning
model rather than to the stopping rule. The model was too cautious under high disturbances with halved time windows,
and the nominal fallback, tested afterwards on new seeds, removed this weakness. In software-in-the-loop flights, the
replanner served fewer tasks late, but the utility differences stayed within the noise. An earlier pre-registered
study, in which only the stopping rule was made execution-aware, found no significant gain in utility
(Section~\ref{sec:study1}).
